\documentclass[11pt,letterpaper]{article}

\usepackage[margin=1in]{geometry}
\usepackage{amsmath,amssymb,amsfonts}
\usepackage{booktabs}
\usepackage{graphicx}
\usepackage{hyperref}
\usepackage{natbib}
\usepackage{enumitem}
\usepackage{xcolor}
\usepackage{float}

\newcommand{\E}{\mathbb{E}}
\newcommand{\Var}{\mathrm{Var}}
\newcommand{\Cov}{\mathrm{Cov}}
\newcommand{\logit}{\mathrm{logit}}
\newcommand{\ATT}{\mathrm{ATT}}

\title{Triply Robust ATT Estimation via Log-Logistic\\Propensity Scores and 2SLS:\\Procedure and Surgical Monte Carlo Results}
\author{}
\date{April 2026 (Updated)}

\begin{document}
\maketitle

\begin{abstract}
We describe a simplified implementation of the triply robust estimator for
the Average Treatment Effect on the Treated (ATT).  The key idea is to
replace the probit balance--propensity-score model~$e_b$ with a
\emph{log-logistic} model whose odds are linear in the covariate:
$e_b(X)/(1-e_b(X)) = \gamma_0 + \gamma_1 f(X)$.  This makes the
augmentation regressor $W = 1/(1-e_b)$ a known linear function of~$f(X)$
after MLE, collapsing the second estimation step from a nonlinear
GMM/ridge-penalized system into a standard two-stage least squares (2SLS)
regression.

A \emph{surgical} Monte Carlo experiment isolates the $e_b$~channel by
using a data-generating process where the outcome regression and $e_a$
are both genuinely misspecified, and the IV regressor space cannot
accidentally span the true untreated mean ($R^2 = 0.62$).  We compare
four estimators:

\begin{itemize}[nosep]
  \item \textbf{S} (simplified): 2SLS with instruments $\{1, q, Z\}$
    --- enforces M3, M4, M5 (no M6).
  \item \textbf{F} (full): 2SLS/GMM with instruments $\{1, q, Z, \hat{o}_b - Z\}$
    --- enforces M3, M4, M5, M6 (overidentified).
  \item \textbf{G} (exactly-identified M6): IV with instruments
    $\{\hat{o}_b - Z, q, Z\}$ --- enforces M4, M5, M6 (no M3).
  \item \textbf{N} (negative control): S with wrong $e_b$ basis.
\end{itemize}

The results show that both~F and~G are consistent under correct~$e_b$
(bias $\to 0$ as $n$ grows), with the G--F gap shrinking to zero.
Estimator~S retains a persistent bias of ${\approx}{-0.025}$ that does
\emph{not} shrink with~$n$, confirming that the M6 moment is
necessary for the $e_b$-channel.

Estimator~G provides a clean resolution: it replaces the
constant-instrument moment M3 with the telescoping moment M6, achieving
exact identification (3 moments, 3 unknowns) while activating the
$e_b$-robust branch.
\end{abstract}


%======================================================================
\section{Introduction and Motivation}
%======================================================================

The triply robust ATT estimator combines three modeling ingredients:
\begin{enumerate}[nosep]
  \item An \emph{outcome regression} (OR) model $m_0(X) = \beta_0 + \beta_1 q(X) + \cdots$,
  \item An \emph{augmentation} propensity score $e_a(X)$ (logit MLE),
  \item A \emph{balance} propensity score $e_b(X)$.
\end{enumerate}
The estimator is consistent for the ATT whenever at least one of these
three components is correctly specified---hence ``triply robust.''

\medskip\noindent
In existing implementations (v1 and v2), $e_b$ is a probit model
$e_b = \Phi(\gamma_0 + \gamma_1 X)$.  Because $1/(1-\Phi(\cdot))$
is nonlinear, recovering the structural parameters
$(\beta_0,\beta_1,\phi)$ from the moment conditions requires a
nonlinear least-squares solver with a ridge penalty~$\lambda$.
This introduces a hyperparameter~$\lambda$ that biases the estimator
at $O(\lambda)$, plus convergence and conditioning issues.

\medskip\noindent
The \textbf{log-logistic simplification} replaces the probit $e_b$ with:
\[
  \frac{e_b(X)}{1 - e_b(X)} \;=\; \gamma_0 + \gamma_1\, f(X),
\]
so that
\[
  W \;\equiv\; \frac{1}{1 - e_b(X)} \;=\; 1 + \gamma_0 + \gamma_1\, f(X),
\]
which is linear in $f(X)$ once $(\gamma_0,\gamma_1)$ are estimated.
This makes the structural estimation a standard \textbf{2SLS} problem.

\medskip\noindent
A crucial question is whether this simplified 2SLS estimator still
inherits the $e_b$-robust consistency branch from the full moment
system.  The original full estimator explicitly enforces an
$e_b$-channel moment (M6) via GMM.  The simplified estimator~S omits
this moment.  Does~S still converge to the correct ATT when only $e_b$
is correctly specified?

This report presents a \emph{surgical} Monte Carlo designed to answer
that question cleanly.


%======================================================================
\section{Estimator: Step-by-Step Procedure}
%======================================================================

\subsection{Setup and Notation}

Let $(Y_i, D_i, X_i)_{i=1}^n$ be an i.i.d.\ sample where $Y_i$ is
the observed outcome, $D_i \in \{0,1\}$ is the treatment indicator,
and $X_i$ is a scalar pre-treatment covariate.  We target the Average
Treatment Effect on the Treated:
\[
  \ATT \;=\; \E[Y(1) - Y(0) \mid D = 1].
\]

\subsection{Estimator S (Simplified)}

\paragraph{Step 1: Logit PS $e_a$ (MLE).}
Fit a logit model $e_a(X) = \logit^{-1}(\alpha_0 + \alpha_1 g(X))$
by standard logit MLE.  Define the fitted odds:
$Z_i = \hat{e}_a(X_i)/(1 - \hat{e}_a(X_i))$.

\paragraph{Step 2: Log-Logistic PS $e_b$ (MLE).}
Fit a log-logistic model:
\[
  \frac{e_b(X)}{1 - e_b(X)} = \gamma_0 + \gamma_1 f(X),
  \qquad \gamma_0 + \gamma_1 f(X_i) > 0 \;\forall\, i.
\]
Define the endogenous regressor $W_i = 1/(1 - \hat{e}_b(X_i)) = 1 + \hat\gamma_0 + \hat\gamma_1 f(X_i)$.

\paragraph{Step 3: 2SLS among controls.}
Restricting to control units ($D_i = 0$), estimate the structural model
\[
  Y_i = \underbrace{\beta_0 + \beta_1 q(X_i)}_{\text{exogenous}} + \phi\, W_i + \varepsilon_i
\]
by 2SLS with instruments $\{1, q(X), Z\}$ and regressors $\{1, q(X), W\}$.

\paragraph{Step 4: ATT computation.}
Impute $\hat{m}_0(X_i) = \hat\beta_0 + \hat\beta_1 q(X_i) + \hat\phi W_i$
for all units.  The ATT estimate is
$\hat\tau_{\ATT} = n_1^{-1}\sum_{i:D_i=1}(Y_i - \hat{m}_0(X_i))$.


\subsection{Estimator F (Full, with M6 Moment)}

Same as~S, except the 2SLS instrument set is expanded from
$\{1, q(X), Z\}$ to $\{1, q(X), Z, \hat{o}_b - Z\}$, where
$\hat{o}_b = \hat{e}_b/(1-\hat{e}_b) = \hat\gamma_0 + \hat\gamma_1 f(X)$
is the fitted odds of~$e_b$.  The additional instrument
$(\hat{o}_b - Z)$ corresponds to the ``telescoping'' M6 moment condition
from the full triply robust moment system.

This makes the 2SLS over-identified (4~instruments for
3~unknowns $(\beta_0, \beta_1, \phi)$).  We use GMM with the efficient
weight matrix $\hat\Omega = (Z'Z)^{-1}$ (2SLS-style).


\subsection{Estimator G (Exactly-Identified, M6 Replaces M3)}

Same as~S, except the constant instrument (1) is \emph{replaced} by
the telescoping instrument $(\hat{o}_b - Z)$:
\begin{itemize}[nosep]
  \item Instruments: $\{(\hat{o}_b - Z),\; q(X),\; Z\}$ (3~instruments).
  \item Regressors: $\{1,\; q(X),\; W\}$ (3~unknowns).
\end{itemize}
This is exactly identified (3 instruments for 3 unknowns), enforcing
moments M4, M5, and M6 exactly.  It drops M3 (the constant-instrument
moment $\E[(1-D)\cdot 1 \cdot R] = 0$).

The key insight is that M3 pins down the intercept~$\beta_0$ to make
the control-sample residuals mean-zero; but when $e_b$ is correct, M6
provides an alternative pinning of the intercept through the
telescoping structure.  By replacing M3 with M6, Estimator~G
explicitly activates the $e_b$-channel at no additional computational
cost and with exact identification.


\subsection{Estimator N (Negative Control)}

Same as~S, but with $e_b$ deliberately misspecified (e.g.,\ linear
odds in~$X$ when the truth is linear odds in~$\exp(X)$).


\subsection{The Quotient-Space Argument}

An important theoretical question is whether~S needs to explicitly
enforce M6 for the $e_b$-robust branch.  The argument for ``no'' is:

\begin{quote}
\itshape
Under the linear-odds specification, the structural nuisance
parameters $(\beta_0, \phi, \gamma_0, \gamma_1)$ are
overparameterized relative to the fitted imputation function.
The ATT estimator depends only on the latter.  Consequently,
under correct specification of the balance propensity score,
consistency does not require that the particular reported nuisance
parameter vector satisfy the telescoping moment condition.  It
suffices that the fitted imputation admit an equivalent structural
representation for which the $e_b$-robust argument applies.
\end{quote}

\noindent
Concretely, the reduced-form imputation is
\[
  m_0(X) = \underbrace{(\beta_0 + \phi(1+\gamma_0))}_{a_0}
           + \beta_1 q(X)
           + \underbrace{(\phi\gamma_1)}_{a_1} f(X).
\]
The ATT depends only on $(a_0, \beta_1, a_1)$, not on the
decomposition into $(\beta_0, \phi, \gamma_0, \gamma_1)$.  If
there exists \emph{some} quadruple
$(\tilde\beta_0, \tilde\phi, \tilde\gamma_0, \tilde\gamma_1)$
with the same $(a_0, \beta_1, a_1)$ that satisfies the
telescoping proof, then~S should give the same limit as~F.

The Monte Carlo below tests whether this argument holds
empirically.

%======================================================================
\section{Surgical Monte Carlo Design}
%======================================================================

\subsection{Design Principles}

To cleanly isolate the $e_b$~channel, the DGP must kill three
escape routes:
\begin{enumerate}[nosep]
  \item The outcome regression must be \emph{truly} wrong.
  \item The $e_a$ model must be \emph{truly} wrong.
  \item The simplified IV regressor space must \emph{not} accidentally
    span the true untreated mean.
\end{enumerate}

\subsection{Data-Generating Process}

\[
  X \sim \mathrm{Uniform}(-2, 2).
\]

\paragraph{True propensity score (linear odds in $e^X$).}
\begin{align}
  \frac{e_0(X)}{1 - e_0(X)} &= 1.5 + 0.8\,e^X, \\
  e_0(X) &= \frac{1.5 + 0.8\,e^X}{2.5 + 0.8\,e^X}.
\end{align}
For $X \in [-2,2]$: propensities range from ${\approx}0.62$ to ${\approx}0.88$.

\paragraph{True outcome model (nonlinear, not in IV span).}
\begin{align}
  \mu_0(X) &= 1 + 0.7\,X^2 + 0.5\,\sin(2X), \\
  Y(0) &= \mu_0(X) + \varepsilon, \quad \varepsilon \sim N(0,1), \\
  Y(1) &= Y(0) + \tau, \quad \tau = 2.
\end{align}

\paragraph{Observed outcome.}
$Y = D\,Y(1) + (1-D)\,Y(0)$, \quad $D \sim \mathrm{Bernoulli}(e_0(X))$.

\paragraph{Why this works.}
\begin{itemize}[nosep]
  \item The log-logistic PS with $f(X)=e^X$ is \emph{correctly specified}.
  \item A wrong $e_a$ model: $\logit^{-1}(\alpha_0 + \alpha_1 X)$
    (logit on $X$, not on $e^X$) is genuinely wrong.
  \item A wrong OR: $r(X;\beta) = \beta_0 + \beta_1 X$ is genuinely wrong.
  \item Most importantly, the true $\mu_0(X) = 1 + 0.7 X^2 + 0.5\sin(2X)$
    is \textbf{not} in the span of $\{1, X, e^X\}$.
\end{itemize}


\subsection{Span Check}

Regressing the true $\mu_0(X)$ on $\{1, X, e^X\}$ on a
large sample ($n = 200{,}000$):
\[
  R^2 = 0.6244.
\]
This confirms that the IV regressor space cannot accidentally span
the truth.  The accidental-specification pathology from the earlier
experiment ($R^2 \approx 1$ when $W$ adds the missing $X$~term)
is eliminated.


\subsection{$e_a$ Misspecification Check}

The logit-on-$X$ approximation has MAE~$= 0.020$ and
correlation~$= 0.96$ with the true PS.  The \emph{functional form}
is wrong (true log-odds are nonlinear in~$X$), but the
\emph{prediction} is reasonable because the PS is nearly monotone
in~$X$ over $[-2,2]$.  This makes the experiment conservative:
any residual bias in~S is despite a ``decent'' $e_a$ approximation.


\subsection{Experimental Cases}

\begin{table}[H]
\centering
\caption{Experimental cases}
\label{tab:cases-surgical}
\begin{tabular}{cllll}
\toprule
\textbf{Case} & \textbf{Estimator} & \textbf{OR} & $e_a$ & $e_b$ \\
\midrule
A-G & G (M6 replaces M3) & $\{1, X\}$ $\times$ & logit($X$) $\times$ & $e^X$ $\checkmark$ \\
A-S & S (simplified, no M6) & $\{1, X\}$ $\times$ & logit($X$) $\times$ & $e^X$ $\checkmark$ \\
A-F & F (full, with M6) & $\{1, X\}$ $\times$ & logit($X$) $\times$ & $e^X$ $\checkmark$ \\
A-N & N (negative control) & $\{1, X\}$ $\times$ & logit($X$) $\times$ & $X$ $\times$ \\
\bottomrule
\end{tabular}
\end{table}

\subsection{Simulation Parameters}

\begin{itemize}[nosep]
  \item Sample sizes: $n \in \{500,\; 2{,}000,\; 10{,}000\}$
  \item Replications: $R = 500$
  \item True ATT: $\tau = 2.0$
  \item Random seed: 42 (reproducible; same draws for all estimators at each~$n$)
\end{itemize}


%======================================================================
\section{Results}
%======================================================================

\subsection{Main Results Table}

\begin{table}[H]
\centering
\caption{Surgical Monte Carlo results ($R=500$).
  \textbf{Bold} marks the key finding: G and F are consistent
  while S has persistent bias.}
\label{tab:results-surgical}
\begin{tabular}{clrrrrrr}
\toprule
\textbf{Case} & \textbf{Estimator} & $n$ &
  \textbf{Mean $\hat\tau$} & \textbf{Bias} & \textbf{SD} &
  \textbf{RMSE} & \textbf{1st-$F$} \\
\midrule
A-G & G (M6 replaces M3) &    500 & 1.9894 & $-$0.0106 & 0.1311 & 0.1315 & 3{,}396 \\
A-S & S (simplified) &    500 & 1.9620 & $-$0.0380 & 0.1404 & 0.1454 & 3{,}396 \\
A-F & F (full+M6)    &    500 & 1.9885 & $-$0.0115 & 0.1340 & 0.1345 & 3{,}396 \\
A-N & N (wrong $e_b$)&    500 & 2.1682 & $+$0.1682 & 0.1482 & 0.2242 &      79 \\
\midrule
A-G & G (M6 replaces M3) & 2{,}000 & 2.0013 & $+$0.0013 & 0.0635 & 0.0635 & 12{,}637 \\
A-S & S (simplified) &  2{,}000 & 1.9746 & $-$0.0254 & 0.0685 & 0.0730 & 12{,}637 \\
A-F & F (full+M6)    &  2{,}000 & 2.0008 & $+$0.0008 & 0.0650 & 0.0650 & 12{,}637 \\
A-N & N (wrong $e_b$)&  2{,}000 & 2.1704 & $+$0.1704 & 0.0720 & 0.1850 &     243 \\
\midrule
\textbf{A-G} & \textbf{G (M6 replaces M3)} & \textbf{10{,}000} &
  \textbf{1.9999} & $\boldsymbol{-0.0001}$ & \textbf{0.0267} &
  \textbf{0.0267} & \textbf{62{,}695} \\
\textbf{A-S} & \textbf{S (simplified)} & \textbf{10{,}000} &
  \textbf{1.9745} & $\boldsymbol{-0.0255}$ & \textbf{0.0288} &
  \textbf{0.0385} & \textbf{62{,}695} \\
\textbf{A-F} & \textbf{F (full+M6)} & \textbf{10{,}000} &
  \textbf{2.0000} & $\boldsymbol{-0.0000}$ & \textbf{0.0273} &
  \textbf{0.0273} & \textbf{62{,}695} \\
A-N & N (wrong $e_b$)& 10{,}000 & 2.1632 & $+$0.1632 & 0.0312 & 0.1661 & 1{,}387 \\
\bottomrule
\end{tabular}
\end{table}


\subsection{G vs.\ F and S vs.\ F Convergence}

\begin{table}[H]
\centering
\caption{Pairwise differences across replications (Case~A, $R=500$).
  G and F converge to the same limit (mean difference~$\to 0$,
  SD shrinking).  S converges to a different limit
  (mean difference constant at ${\approx}{-0.026}$).}
\label{tab:sf-diff}
\begin{tabular}{rrrrr}
\toprule
$n$ & Mean($\hat\tau_G - \hat\tau_F$) &
  SD($\hat\tau_G - \hat\tau_F$) &
  Mean($\hat\tau_S - \hat\tau_F$) &
  Mean($\hat\tau_G - \hat\tau_S$) \\
\midrule
     500 & $+$0.0009 & 0.0099 & $-$0.0265 & $+$0.0274 \\
   2{,}000 & $+$0.0005 & 0.0046 & $-$0.0262 & $+$0.0266 \\
  10{,}000 & $-$0.0000 & 0.0019 & $-$0.0255 & $+$0.0255 \\
\bottomrule
\end{tabular}
\end{table}

\begin{figure}[H]
\centering
\includegraphics[width=0.85\textwidth]{centered_monte_carlo.png}
\caption{Left: Bias as a function of $n$ for all four estimators.
  G and F converge to zero; S plateaus at ${\approx}{-0.025}$;
  N at ${\approx}{+0.17}$.
  Right: Distribution of $\hat\tau_G - \hat\tau_F$ at each sample size
  (concentrating around zero with shrinking spread).}
\label{fig:gf-diff}
\end{figure}


%======================================================================
\section{Interpretation}
%======================================================================

\subsection{The Central Finding}

The experiment cleanly separates the estimators under a DGP where
\emph{only} $e_b$ is correctly specified:

\begin{enumerate}
  \item \textbf{Estimator~G (M6 replaces M3) is consistent.}
    Bias at $n = 10{,}000$ is $-0.0001$, essentially zero.  The bias
    collapses: $-0.011 \to +0.001 \to -0.000$ as $n$ grows.

  \item \textbf{Estimator~F (full, overidentified with M6) is consistent.}
    Bias at $n = 10{,}000$ is $-0.0000$.  The bias
    shrinks monotonically: $-0.012 \to +0.001 \to -0.000$.

  \item \textbf{G and F converge to the same limit.}
    The mean $\hat\tau_G - \hat\tau_F$ is $+0.0009 \to +0.0005 \to
    -0.0000$ and the SD shrinks from $0.010$ to $0.002$.  This
    confirms they are asymptotically equivalent.

  \item \textbf{Estimator~S (simplified, without M6) has a persistent bias.}
    Bias at $n = 10{,}000$ is $-0.026$, essentially unchanged from
    $n = 2{,}000$ ($-0.025$).  S converges to a different
    probability limit from G and~F.

  \item \textbf{Estimator~N (wrong $e_b$) has clear, large bias.}
    Stable at $+0.17$ across all~$n$, confirming it is a valid negative
    control.
\end{enumerate}


\subsection{Why S Fails and G Succeeds}

The quotient-space argument (Section~2.4) says: if the reduced-form
imputation $(a_0, \beta_1, a_1)$ can be ``re-decomposed'' into a
structural quadruple that satisfies the telescoping proof, then~S
should be as good as~F.

The Monte Carlo shows this is \emph{not sufficient} for S: while
both S and~F produce imputations in the same reduced-form
space $\{1, X, e^X\}$, they use \emph{different instrument sets}.
S uses instruments $\{1, X, \hat{o}_a\}$ (enforcing M3, M4, M5);
F and G use the telescoping instrument $(\hat{o}_b - \hat{o}_a)$
to enforce M6.  When the OR is misspecified, the projection of
$\mu_0$ onto $\{1, X, e^X\}$ depends on the instruments.  Only
the projection selected by the M6 instrument gives the correct
ATT-weighted average.

\medskip\noindent
Estimator~G cleanly resolves this by \emph{replacing} M3 with M6:
\begin{itemize}[nosep]
  \item S uses $\{1, q, Z\}$ for $\{1, q, W\}$ $\to$ enforces M3, M4, M5.
  \item G uses $\{(\hat{o}_b - Z), q, Z\}$ for $\{1, q, W\}$
    $\to$ enforces M4, M5, M6.
  \item F uses $\{1, q, Z, (\hat{o}_b - Z)\}$ for $\{1, q, W\}$
    $\to$ enforces M3, M4, M5, M6 (overidentified).
\end{itemize}
G is exactly identified (3 moments, 3 unknowns) and achieves the
same consistency as the overidentified~F, while dropping the
constant-instrument moment that S uses instead.  This pinpoints
M6 as the critical moment: it is both necessary and sufficient
(together with M4, M5) for the $e_b$-channel.


\subsection{Estimator F's Consistency Mechanism}

To see why F works, note that its instrument set
$\{1, X, \hat{o}_a, \hat{o}_b - \hat{o}_a\}$ can be rewritten as
$\{1, X, \hat{o}_a, \hat{o}_b\}$.  The additional instrument
$\hat{o}_b = \hat\gamma_0 + \hat\gamma_1 e^X$ is a function of
$e^X$ alone (up to constants).  When $e_b$ is correct, the
telescoping structure ensures that the population moment
\[
  \E\bigl[(1-D)\,(\hat{o}_b - \hat{o}_a)\,(Y - m_0(X))\bigr] = 0
\]
holds at the true $m_0$, even when $m_0$ is not in the model's
regressor space.  This gives F an additional correctly-specified
moment condition that identifies the right projection.

S, by contrast, relies only on $\{1, X, \hat{o}_a\}$.  When OR
is wrong, $\hat{o}_a$ (which is a function of the wrong $e_a$)
does not provide the same identifying information as
$(\hat{o}_b - \hat{o}_a)$.  Hence~S converges to the wrong
projection.


\subsection{Practical Implications}

\begin{enumerate}[nosep]
  \item \textbf{The simplified 2SLS estimator~S is NOT triply robust.}
    It retains the OR-robustness and $e_a$-robustness branches, but the
    $e_b$-channel requires explicit activation via M6.

  \item \textbf{Estimator~G provides the cleanest fix.}
    Replacing the constant instrument with $(\hat{o}_b - Z)$ gives
    exact identification (3 moments, 3 unknowns) while activating the
    $e_b$-channel.  No overidentification, no GMM weight matrix.

  \item \textbf{Including M6 is computationally cheap.}
    Both G and~F add only one column to the instrument matrix.
    The computation remains a single linear system (one matrix solve),
    with no nonlinear solver or ridge penalty.

  \item \textbf{G and F are asymptotically equivalent.}
    The G--F gap shrinks to zero, confirming they target the same
    probability limit.  G is simpler (exactly identified);
    F may have slightly better finite-sample efficiency from
    overidentification.

  \item \textbf{Recommendation: use~G or~F.}
    Either estimator restores the $e_b$-channel consistency.
    G is preferred for simplicity (exact identification);
    F is preferred if one wants to exploit M3 for additional efficiency.
\end{enumerate}


%======================================================================
\section{Comparison with Earlier 8-Case Experiment}
%======================================================================

The earlier Monte Carlo (Appendix) used 8~cases across two DGPs
with $X \sim \mathrm{Uniform}(10,20)$.  Cases~5--7 (DGP~2, $e_b$
correct) appeared to show that the simplified estimator was
consistent under correct~$e_b$.  However, those DGPs used
$f(X) = X$ and $\mu_0(X) = 2 + 0.5\log X + 0.3 X$, so the
IV regressor space $\{1, \log X, X\}$ \emph{accidentally spanned}
the true $\mu_0$---the accidental-specification pathology flagged
in Section~4 of the previous report.

The surgical design eliminates this by using $f(X) = e^X$ and a
$\mu_0$ that includes $X^2$ and $\sin(2X)$, giving $R^2 = 0.62$.
Under this clean design, the simplified estimator's $e_b$~channel
fails, confirming that the earlier apparent success was due to
accidental spanning, not genuine $e_b$-channel robustness.


%======================================================================
\section{Implementation Summary}
%======================================================================

The implementation consists of three Python files:

\paragraph{\texttt{triple\_robust\_iv.py}} contains:
\begin{itemize}[nosep]
  \item \texttt{triply\_robust\_iv(...)}: Estimator~S (simplified 2SLS).
  \item \texttt{triply\_robust\_iv\_full(...)}: Estimator~F
    (full 2SLS/GMM with M6 instrument).
  \item \texttt{triply\_robust\_iv\_centered(...)}: Estimator~G
    (exactly-identified, M6 replaces M3).
  \item Shared helpers: \texttt{\_log\_logistic\_mle()},
    \texttt{\_invlogit()}, \texttt{\_clip\_ps()}.
  \item Supports $e_b$ bases: $X$, $\sin X$, $e^X$.
  \item Supports OR columns: $\{1,\log X\}$, $\{1,\log X, X\}$, $\{1, X\}$.
\end{itemize}

\paragraph{\texttt{simulation\_centered.py}} contains:
\begin{itemize}[nosep]
  \item \texttt{dgp\_surgical(rng, n, kappa)}: the surgical DGP.
  \item \texttt{run\_diagnostics()}: span check ($R^2$) and
    $e_a$ misspecification check.
  \item \texttt{run\_simulation(R, sample\_sizes)}: Monte Carlo loop
    running G, S, F, and N at each~$n$.
\end{itemize}

\paragraph{\texttt{simulation\_surgical.py}} contains the earlier S-vs-F
experiment (kept for reference).

\paragraph{\texttt{simulation\_iv.py}} contains the earlier 8-case
experiment (kept for reference).

\paragraph{Usage:}
\begin{verbatim}
  python simulation_centered.py             # R=10  smoke test
  python simulation_centered.py 500         # R=500, as in this report
\end{verbatim}


%======================================================================
\section{Conclusion}
%======================================================================

\begin{enumerate}
  \item The log-logistic reparameterization of~$e_b$ successfully
    reduces the imputation step to 2SLS.  No nonlinear solver,
    ridge penalty, or convergence monitoring is needed.

  \item The simplified 2SLS (Estimator~S) does \textbf{not}
    inherit the $e_b$-robust consistency branch.  Under correct $e_b$
    with misspecified OR and $e_a$, S has a persistent bias of
    ${\approx}{-0.025}$ that does not shrink with~$n$.

  \item \textbf{Estimator~G}, which replaces the constant instrument
    (M3) with the telescoping instrument (M6), is exactly identified
    and \textbf{is} consistent under correct $e_b$.  Its bias
    collapses to zero as $n$ grows, matching the overidentified~F.

  \item The \textbf{full} estimator~F (overidentified, M3+M4+M5+M6)
    is also consistent.  G and~F converge to the same probability
    limit, with the G--F gap shrinking at rate $O(1/\sqrt{n})$.

  \item The quotient-space argument---that many structural parameter
    vectors generate the same fitted imputation, so M6 need not be
    explicitly enforced---is \textbf{not sufficient} for S's consistency.
    The M6 moment is genuinely needed to select the right projection
    of $\mu_0$ onto the model space.

  \item \textbf{Recommendation:} use Estimator~G or~F.  G is the
    simpler choice (exact identification, single matrix solve);
    F may offer modest finite-sample efficiency gains from
    overidentification.  Both retain all three robustness channels.
\end{enumerate}


%======================================================================
\appendix
\section{Earlier 8-Case Experiment (Reference)}
\label{app:eight-case}
%======================================================================

For completeness, the earlier experiment is summarised below.  Two DGPs
with $X \sim \mathrm{Uniform}(10,20)$ and true $\ATT = 3$:

\paragraph{DGP~1 (Logit-in-sin PS).}
$e_0 = \logit^{-1}(-1 + 0.1\sin X)$,\;
$\mu_0 = 2 + 0.5\log X + 0.3 X$,\;
$\tau = 3$.

\paragraph{DGP~2 (Linear-odds PS).}
$\mathrm{odds} = -0.5 + 0.05 X$,\;
$\mu_0 = 2 + 0.5\log X + 0.3 X$ (same),\;
$\tau = 3$.

\begin{table}[H]
\centering
\caption{Earlier 8-case results ($R=200$, $n=5{,}000$, Estimator~S only).
All cases show low bias, but Cases~5--7 (DGP~2, $e_b$ correct) are
contaminated by accidental spanning ($R^2 \approx 1$).}
\label{tab:results-old}
\begin{tabular}{ccccccccc}
\toprule
\textbf{Case} & \textbf{DGP} & \textbf{OR} & $e_a$ & $e_b$ & \textbf{Mean $\hat\tau$} & \textbf{Bias} & \textbf{SD} & \textbf{RMSE} \\
\midrule
1  & 1 & $\checkmark$ & $\checkmark$ & $\times$ & 2.999 & $-$0.001 & 0.033 & 0.033 \\
2  & 1 & $\checkmark$ & $\times$ & $\times$ & 3.003 & $+$0.003 & 0.032 & 0.032 \\
3  & 1 & $\times$ & $\checkmark$ & $\times$ & 2.999 & $-$0.002 & 0.033 & 0.033 \\
4  & 1 & $\times$ & $\times$ & $\times$ & 3.000 & $+$0.000 & 0.032 & 0.032 \\
\midrule
5  & 2 & $\times$ & $\times$ & $\checkmark$ & 3.005 & $+$0.005 & 0.038 & 0.038 \\
6  & 2 & $\checkmark$ & $\times$ & $\checkmark$ & 2.999 & $-$0.001 & 0.038 & 0.038 \\
7  & 2 & $\times$ & $\times$ & $\checkmark$ & 3.001 & $+$0.001 & 0.036 & 0.036 \\
4b & 2 & $\times$ & $\times$ & $\times$ & 2.996 & $-$0.004 & 0.041 & 0.041 \\
\bottomrule
\end{tabular}
\end{table}


\end{document}
